
\frontslide{Gestion des mémoires}

\begin{frame}[fragile]{Gestion des mémoires}

Contenu de ce cours\,:

\begin{redbullet}
   \item Mémoires gérées par un SGBD: mémoire cache, disque.

  \item La gestion des lectures.

  \item La gestion des écritures.
  
  \item Quelques applications du principe de localité.
\end{redbullet}


\vfill
\begin{blocgauche}
\alert{Ces diapositives correspondent au support en ligne disponible sur
le site \url{http://sys.bdpedia.fr/stock.html\#s2-gestion-des-memoires}}
\end{blocgauche}

\end{frame}


\begin{frame}{Le cache et le disque}

Paramétrage des SGBD: partie de la mémoire centrale qui sert de {\it cache}.

\vfill

\figSlideWithSize{buffer-disque}{7cm}


\end{frame}


\begin{frame}{Opération de lecture}

Toute opération de lecture contient l'adresse du bloc à lire. 

\begin{redbullet}
   \item \alert{Si la donnée fait partie d'un bloc dans le cache}, 
      le SGBD prend le bloc, accède à la donnée dans le bloc,  et la retourne;
   \item  \alert{Sinon} il faut d'abord lire un bloc du disque, et le placer dans le cache, 
   pour se ramener au cas précédent.
\end{redbullet}

\vfill

\begin{blocgauche}
\alert{Hypothèse de localité}: le SGBD garde en mémoire
les blocs après utilisation, afin d'exploiter 
à la fois 

\begin{redbullet}
  \item la \alert{localité spatiale}: les autres données du bloc.
   \item la \alert{localité temporelle}: il est probable qu'on va relire la donnée dans peu de temps.
\end{redbullet}
\end{blocgauche}
\end{frame}

\begin{frame}{Lectures logiques, lectures physiques\,: le \textit{Hit ratio}}

Le paramètre
qui mesure l'efficacité d'une mémoire cache
est le \textit{hit ratio}\,:
\begin{center}
   $hit\ ratio\ =\ \frac{nb\ de\ lectures\ logiques - nb\ lectures\ physiques}
         {nb\ de\ lectures\ logiques}$
\end{center}
\begin{redbullet}
 \item Si toutes les lectures logiques (demande de bloc) aboutissent
à une lecture physique (accès au disque), le \textit{hit
ratio} est 0.

\item Aucune lecture physique (tout est en mémoire):
le \textit{hit ratio} est de 1.
\end{redbullet}

\vfill
\begin{blocgauche}
\alert{En pratique}. Un bon \textit{Hit Ratio} est supérieur à 80\%, voire 90\%:
presque toutes les lectures se font en mémoire!
\end{blocgauche}

\end{frame}

\begin{frame}{Comment avoir un bon \textit{Hit ratio}?}


Evident: 

\vfill

\begin{redbullet}
  \item il faut allouer le plus de mémoire possible au SGBD.
  \item il faut limiter la taille de la base (attention aux longs textes, aux images, aux données binaires...)
\end{redbullet}

\vfill

\begin{blocgauche}

\alert{Mais surtout: il faut que les données \textbf{utiles} soient
en mémoire centrale}

\medskip

Intuition: certaines parties de la base sont \alert{beaucoup} plus lues
que d'autres: il faut qu'elles tiennent en mémoire.
\end{blocgauche}

\end{frame}

\begin{frame}{Stratégie de remplacement}


Que faire quand la mémoire est pleine?

\vfill

On applique le principe de localité: l'algorithme le plus courant est dit \textit{Least Recently Used} (LRU).

\begin{redbullet}
  \item  La "victime" est le bloc dont la dernière date de lecture logique est la plus ancienne. 
  \item Ce bloc est alors soustrait de la mémoire centrale 
  \item  Le nouveau bloc vient le remplacer.
\end{redbullet}

\vfill

\begin{blocgauche}
\alert{Conséquence}: le contenu du cache est une image fidèle de l'activité \alert{récente} sur la base de données. 
\end{blocgauche}
\end{frame}


\begin{frame}{Opération de mise à jour}

Approche naïve: on trouve le bloc (comme pour une lecture), on modifie la donnée.

\vfill
 
\figSlideWithSize{write-naif}{7.5cm}

\vfill
\begin{blocgauche}
Faut-il écrire le bloc ou non? Deux mauvais choix (pourquoi?)
\end{blocgauche}
\end{frame}


\begin{frame}{La bonne méthode; le fichier journal}

On gère la base (accès alléatoire) \alert{et} un fichier séquentiel, le \alert{journal} (log).

\medskip
 
\figSlideWithSize{recovery}{7.5cm}

\medskip
\begin{blocgauche}
Problème résolu: données sur disque (sûres) et écritures séquentielles.
\end{blocgauche}
\end{frame}

\begin{frame}{Le principe de localité et ses conséquences}


\alert{Principe de localité}\,: l'ensemble des  données utilisées
par une application pendant une période donnée forme souvent
un groupe bien identifié.


\begin{redbullet}
 \item \alert{Localité spatiale}\,: si une donnée $d$ est utilisée,
    les données proches de $d$ ont de fortes chances de l'être
    également.

 \item \alert{Localité temporelle}\,: quand une application accède
   à une donnée $d$, il y a de fortes chances qu'elle y accède à nouveau
   peu de temps après.

 \item \alert{Localité de référence}\,: si une donnée $d_1$ référence
     une donnée $d_2$, l'accès à $d_1$ entraîne souvent
    l'accès à $d_2$.
\end{redbullet}

\begin{blocgauche}
Les systèmes exploitent ce principe 
en déplaçant dans la hiérarchie des mémoires des groupes
de données.
\end{blocgauche}

\end{frame}

\begin{frame}{Résumé}

Un critère essentiel de la performance est la taille de la mémoire cache.

\vfill 

Pour vérifier que le système est bien paramétré: on calcule le \textit{Hit Ratio}. Il doit être supérieur à 80\%

\vfill

Sinon, les pistes:

\vfill
\begin{blocgauche}
\begin{redbullet}
  \item Vérifier qu'on ne lit que des données utiles.
  \item Regarder si on peut limiter la taille des données (ex: partitionner la table
       en mettant les colonnes volumineuses à part).
   \item Ajouter de la mémoire...
\end{redbullet}
\end{blocgauche}

\end{frame}
